\section{Summary and conclusions}\label{sec:conclusions}
\sectionstate{partial}{new 2026-11-10: a draft written from the results of this version. The table collects numbers that are
each in the body (sources per row: \texttt{papers/paper1\_v2/NUMBERS\_RUNDOWN.md}, item 15). The paragraph on the design choices
holds the results that the old section 2.1 quoted before the detector was described; ``what could be recovered'' is the old
section 5.6 in three sentences (the long version is in \cref{app:moved}).}
\intent{What the chambers do, what varies between them, what the design choices bought, and what limits the result.}

\Cref{tab:summary} collects the performance of the five chambers.

\begin{table}[tbp]
  \centering\small
  \caption{The five chambers on the cosmic bench (chamber~E's position resolution is from the beam). Efficiencies are
  kept-track efficiencies in the interior (\cref{sec:effdef}); chamber~A's single-track efficiency over the full active
  area is 94.8\,\%. C, D and E were calibrated on the runs they are quoted on. Dashes: not measured.}
  \label{tab:summary}
  \setlength{\tabcolsep}{4.5pt}
  \begin{tabular}{lccccc}
    \hline
    & A & B & C & D & E \\
    \hline
    long run, resist / drift HV [V] & 490 / 1000 & 495 / 1000 & 495 / 700 & 495 / 700 & 495 / 600 \\
    efficiency at the long-run voltage [\%] & 92.3 & 90.3 & 74.4 & 68.0 & 40.9 \\
    efficiency at the best voltage [\%] & 93.8 & 92.2 & 85.8 & 89.8 & -- \\
    best voltage [V] & 470 & 480 & 480 & 430 & -- \\
    angular resolution X / Y [\textdegree] & 0.99 / 0.90 & 1.05 / 1.13 & 1.26 / 1.51 & 1.23 / 1.46 & 2.41 / 2.08 \\
    position resolution X / Y [\textmu m] & 128 / 113 & 163 / 176 & -- & -- & 138 / 120 \\
    event time; uncorrelated part [ns] & 41; 22 & 42; 28 & 78; 53 & 46; 26 & -- \\
    X fraction of the charge & 0.485 & 0.513 & 0.504 & 0.485 & 0.509 \\
    dish of the board at 150\,mm [mm] & 1.31 & 1.87 & 1.76 & 1.68 & 1.69 \\
    water in the gas [\%] & 0.97 & 0.69 & 1.00 & 0.48 & 0.42 \\
    dark current at 490\,V [\textmu A] & 0.3--1.1 & 1.0--2.4 & 0.45 & 6.0--6.4 & 0.58 \\
    \hline
  \end{tabular}
\end{table}

\paragraph{What a chamber delivers.} A compact resistive Micromegas TPC with a 28\,mm drift gap measures the direction
of a track from a single chamber to 1.0\textdegree{} and 0.9\textdegree{} in its two views, independently of the angle,
the gain and the drift field, and out of sample. That is the quantity the experiment needs
(\cref{sec:requirement}), and it is below the multiple scattering that leptons of 3 to 16\,MeV undergo in the air and
in the entrance foils before they are measured. The chamber needs no external calibration of its drift: two velocities
and a geometric gap, all from the data, agree to 2\,\%, and the velocity doubles as a monitor of the water in the gas.
It detects 94.8\,\% of single tracks over its full area, pairs the two views of two-particle events correctly in
97.7\,\%, reconstructs both of two tracks more than 18\,mm apart in 95\,\% of cases with their opening angle known to
1.3\textdegree, resolves positions to about 120\,\textmu m and times to 40\,ns.

\paragraph{What varies between chambers.} Two of five chambers, A and B, reach this. Chambers C and D reach 86 and
90\,\% efficiency at their best voltage but have no plateau: discharges set in before it, and at the voltage of the
long runs they cost 17 and 24\,\% of the tracks. Chamber~D also has a fault on half of its X readout, and 62\,\% of the area
of chamber~E does not amplify, in stripes (\cref{app:beam}). What does not vary is the geometry: all five boards are dished by 1.3 to 1.9\,mm, and all five share the charge equally between
the views.
\askStephan{What distinguishes boards A and B from C, D and E in production? This paragraph is where the paper should say
it.}
\askDylan{``62\,\% of the area of chamber~E does not amplify'' is the old number; the story file replaced it by ``52\,\% of the $x$
range below half the plateau; 41\,\% two-view detection''. Which do you want?}

\paragraph{The design choices, judged by the data.} \emph{Pads for an equal charge share} worked: the X strips read
0.485 to 0.513 of the charge, with a track-to-track spread of 6\,\% in the ratio, and that equality is what pairs the
views (\cref{sec:pairing}). \emph{The resistive layer} did what it was chosen for: no dead time is measurable after a
discharge. Its cost is in the reconstruction, which has to model the spreading of the charge (\cref{sec:reco}).
\emph{Resistive strips} draw a dark current of about 1\,\textmu A at the working point; no board with a uniform layer
was on the bench to compare. \emph{The orientation of the strips} put the 10\,\% of area lost to passivation where it
was wanted, but the better resolution expected across the strips is not seen: chamber~A measures the angle better
along them (0.99\textdegree{} across, 0.90\textdegree{} along), chamber~B the reverse, and the position residuals of the
two views agree to 12\,\textmu m. The two sides left without passivation fringe over 20 to 30\,mm, which costs a
fiducial margin for the angle and 0.9 points of efficiency.

\paragraph{What limits it, and what could be recovered.} The chamber hardware accounts for 0.6 points of
inefficiency; it alone would allow 99.4\,\%. The rest is real ionisation that a one-straight-track reconstruction does
not describe: $\delta$-rays (2.4 points), the edge zone (0.9), a second particle merged in one projection. A
reconstruction that recognises the topology before it fits, a drift model that knows the edges and a veto that counts
strips per cluster are the named routes, and none of them has been measured. For the leptons the chambers were built
for, the limit is elsewhere: scattering in the gas while the track is measured, and in the foils and the air before
(\cref{sec:straight,sec:pointing}).

\tbw{What the n\_TOF campaign used the chambers for; comparison with other micro-TPC and resistive Micromegas results in the
literature (ATLAS NSW micro-TPC, T2K resistive Micromegas TPCs, charge-dispersion readouts): one paragraph, with numbers.}
